\documentclass[conference]{IEEEtran}
\IEEEoverridecommandlockouts
% The preceding line is only needed to identify funding in the first footnote. If that is unneeded, please comment it out.
\usepackage{cite}
\usepackage{amsmath,amssymb,amsfonts}
\usepackage{algorithmic}
\usepackage{graphicx}
\usepackage{textcomp}
\usepackage{xcolor}
\def\BibTeX{{\rm B\kern-.05em{\sc i\kern-.025em b}\kern-.08em
    T\kern-.1667em\lower.7ex\hbox{E}\kern-.125emX}}
\begin{document}

\title{AprilTag Follower\\
{\footnotesize Group \#1, Construction Robotics, WS 2022/23}
}

\author{Utku Elagoz (532909), Moayed Mohamed Hani Ibrahim Elmoursi (Student ID), Rami Nasr (Student ID),\\ Ararat Cetinkaya (Student ID), and Abdelrahman Yousef Elhossiny Abounida (Student ID)}

\maketitle
% add page numbers
\thispagestyle{plain}
\pagestyle{plain}

\begin{abstract}
ADD Abstract Here
\end{abstract}

\begin{IEEEkeywords}
keyword 1, keyword 2, keyword 3
\end{IEEEkeywords}

\section{Introduction}
ADD INTRODUCTION HERE
% \textit{Scientific Writing in Engineering} book \cite{b1} and the lecture slides.

\section{Methodology}
    \subsection{Camera Calibration}
        Cameras are an important part of many different areas in robotics such as mobile robots, robot dogs and medical robots. They are widely used for navigation, localization, mapping, collision avoidance, object detection, and object or marker following. A camera basically converts the light particles which are hit the camera’s image sensor to digital images which consist of pixels, which are the smallest units of digital images. There are various camera options used in robotics, nevertheless, none of those cameras are perfect to transform the real world into digital images \cite{b1}. The physical flaws contained by cameras can lead to distortion of the images. There exist different forms of distortions. Two significant ones are radial distortion, where straight lines turn curved and tangential distortion, where the actual scene looks stretched or tilted in the image \cite{b2}. The visual examples of the distortion types can be seen in \\Figure \ref{fig:fig1}.
        
        \begin{figure}[h]
            \centering
            \includegraphics[width=200]{"Example-of-the-distortion-created-by-the-lens-of-a-camera"}
            \caption{Undistorted, Radial distortion,  Tangential Distortion \cite{b1}}
            \label{fig:fig1}
        \end{figure}
        
            
        For many robotics applications that use cameras, precision and accuracy are key elements, as a result, correcting these distortions is critical. Camera calibration is the process of correcting the distortions by estimating the camera parameters called intrinsic and extrinsic. Extrinsic parameters point to the orientation and location of the camera frame with respect to the real-world frame. On the other hand, intrinsic parameters point to the internal parameters of the camera and the relationship between camera coordinates and coordinates of the pixels of the digital image in the image frame. The extrinsic parameters are rotation (3x3) and translation (3x1) matrices, whereas the intrinsic parameters are focal length, skew factor, the optical center of the image, etc. The intrinsic parameters can be shown in a Calibration Matrix, K where f is the focal length and c is the center pixel \cite{b1}.\\
        
        K =
        \begin{bmatrix}
            f_x & s & c_x\\
            0 & f_y & c_y\\
            0 & 0 & 1
        \end{bmatrix}\\
    
        Chessboard calibration is a common method used in camera calibration due to features such as not having the z-axis because of being a flat surface and having easily detectable corners and points which appear at straight lines using the Hough transform or other methods. There are different approaches to camera calibration, the most commonly used ones can be named Zhang’s Method, Tsai’s method, and Direct Linear Transformation. Zhang’s method uses a checkerboard pattern in which is placed in at least more than two different positions. After extracting the corners of the pattern, it calculates a projective transformation between the captured image points unto a scaling factor. Later on, the intrinsic and extrinsic parameters are solved. On the other hand, Tsai’s method has an assumption that if (some) camera parameters are provided, which reduces the initial estimation guesses \cite{b1}\cite{b2}\cite{b3}.\\
        In the project, OpenCV library with python is used to perform camera calibration. 20 pictures of a checkerboard with different poses are collected for camera calibration and stored in a file. Each image is grey scaled and corners of the checkerboards are found using OpenCV’s findChessboardCorners function. After, using calibrateCamera function which uses Zhang’s method, the calibration matrix for intrinsic parameters and the distortion coefficients are calculated.
    
    
    \subsection{AprilTag Detection}
        To detect the April tags in the images, Apriltag and OpenCV libraries are used. Functions from OpenCV are utilized for the reading and greyscaling the images. Later, the detect function of the Apriltag library is used for detecting the position of the Apriltags in the images. 

        The detect function consists of two main algorithms which are fast corner detection and quadrilateral detection. Fast corner detection algorithms are used to quickly identify potential corners of Apriltags in the image, which are then used as starting points for quadrilateral detection. Once potential corners have been identified, quadrilateral detection algorithms are used to search for quadrilaterals which are basically four-sided shapes in the image that match the expected shape of AprilTag markers.

        Once a marker is detected, the marker's binary payload is decoded to determine the marker's identifier, and the position and orientation of the marker are estimated using methods such as perspective-n-point (PnP) algorithms. Finally, the detected marker is validated using various checks such as verifying the parity of the marker's identifier and verifying that the marker's quadrilateral has the correct shape.

        Before starting to apply these algorithms, detect function preprocess the images by modifying them to reduce noise and enhance contrast. Which helps the fast corner detection algorithm to operate smoothly. 
    
    \subsection{Controlling of the mIDOG }

Keep your text and graphic files separate until after the text has been 
formatted and styled. Do not number text heads---{\LaTeX} will do that 
for you.

\section{Results}

\section{Discussion}

\section{Conclusion}

% \section{References}

\begin{thebibliography}{00}
\bibitem{b1} K. Smarsly and K. Dragos, Scientific Writing the in Engineering. Hamburg, Germany: tredition, 2019.
\bibitem{b2} G. Eason, B. Noble, and I. N. Sneddon, ``On certain integrals of Lipschitz-Hankel type involving products of Bessel functions,'' Phil. Trans. Roy. Soc. London, vol. A247, pp. 529--551, April 1955.
\bibitem{b3} J. Clerk Maxwell, A Treatise on Electricity and Magnetism, 3rd ed., vol. 2. Oxford: Clarendon, 1892, pp.68--73.
\bibitem{b4} I. S. Jacobs and C. P. Bean, ``Fine particles, thithe n films and exchange anisotropy,'' in Magnetism, vol. III, G. T. Rado and H. Suhl, Eds. New York: Academic, 1963, pp. 271--350.
\bibitem{b5} K. Elissa, ``Title of paper if known,'' unpublished.
\bibitem{b6} R. Nicole, ``Title of paper with only first word capitalized,'' J. Name Stand. Abbrev., in press.
\bibitem{b7} Y. Yorozu, M. Hirano, K. Oka, and Y. Tagawa, ``Electron spectroscopy studies on magneto-optical media and plastic substrate interface,'' IEEE Transl. J. Magn. Japan, vol. 2, pp. 740--741, August 1987 [Digests 9th Annual Conf. Magnetics Japan, p. 301, 1982].
\bibitem{b8} M. Young, The Technical Writer's Handbook. Mill Valley, CA: University Science, 1989.
\end{thebibliography}
\vspace{12pt}

\end{document}
